% Generated draft: requires mathematical and source review.
\documentclass[11pt]{article}
\usepackage[margin=25mm]{geometry}
\usepackage{fontspec}
\setmainfont{Latin Modern Roman}
\usepackage{amsmath,amssymb,mathtools,unicode-math}
\setmathfont{Latin Modern Math}
\usepackage{enumitem,needspace,xurl,hyperref}
\setlength{\parindent}{0pt}
\setlength{\parskip}{5pt}
\setlength{\emergencystretch}{3em}
\newcommand{\N}{\mathbb{N}}
\newcommand{\Z}{\mathbb{Z}}
\newcommand{\Q}{\mathbb{Q}}
\newcommand{\R}{\mathbb{R}}
\newcommand{\C}{\mathbb{C}}
\newcommand{\F}{\mathbb{F}}
\newcommand{\A}{\mathbb{A}}
\newcommand{\PP}{\mathbb{P}}
\newcommand{\E}{\mathbb{E}}
\newcommand{\eps}{\varepsilon}
\newcommand{\dd}{\,\mathrm{d}}
\begin{document}
\section*{UnsolvedMath ID: 30003935}
\textbf{Source problem code:} catalogue entry
\subsection{Definitions and mathematical statement}
\textbf{Source question.} Does the discrete ensemble Kalman inversion iteration converge strongly to continuous paths of its limiting stochastic system in the general nonlinear setting?

Here ensemble Kalman inversion (EKI) denotes the method named in the supplied record. The record supplies neither its discrete recurrence nor the limiting stochastic equation. Consequently, it does not specify the mathematical objects whose convergence is requested.

To describe the missing convergence specification, let \(h>0\) denote a prospective time step, \(U^h\) a prospective interpolation of the discrete ensemble, and \(U\) the prospective continuous limiting ensemble process. These are placeholders, not definitions of a particular EKI scheme. For example, if both processes were defined on a common probability space with values in a specified normed ensemble space \(E\), one possible meaning of strong convergence would be
\[
\lim_{h\downarrow0}\mathbb{E}\!\left[\sup_{0\leq t\leq T}\|U^h(t)-U(t)\|_E^p\right]=0.
\]
This expression is illustrative only: the source does not specify an exponent \(p\), a time horizon \(T\), a norm, a coupling, or whether the supremum belongs inside the expectation. Convergence in moments at each fixed time and other notions of strong convergence would give different statements.

\textbf{Unresolved formulation.} Determine whether the source's discrete stochastic EKI converges strongly to its limiting continuous stochastic system for the intended class of nonlinear forward maps. A precise quantified statement requires the discrete update, limiting equation, ensemble and observation spaces, noise and initial-data assumptions, interpolation and coupling, admissible nonlinear maps, and intended convergence topology. None of these is specified in the supplied record, so no particular completion is asserted here.
\subsection{Short English statement}
The question asks whether stochastic ensemble Kalman inversion approximates its continuous stochastic limit strongly for general nonlinear forward maps. The supplied record does not define the iteration, limiting equation, hypotheses, or meaning of strong convergence, so the intended mathematical statement remains incomplete. Its partially solved label and literature assessment are supplied dated curation, not independently verified conclusions.
\subsection{Sources}
\begin{itemize}
\item Supplied local project registry record: permanent UnsolvedMath ID \(30003935\), title ``Strong Convergence of Ensemble Kalman Inversion,'' source problem code \texttt{catalogue entry}, collection ``Oberwolfach Reports Open Problems.'' The record gives proposed year 2018 and names no proposer.
\item Supplied report citation: \emph{New Directions in Stochastic Optimisation} (2019), Oberwolfach Reports, pp.\ 2303--2384, DOI \texttt{10.4171/owr/2018/38}. \url{https://doi.org/10.4171/owr/2018/38}. The supplied background describes these reports as documentation of mathematical workshops and their open-problem discussions. The report contents have not been independently inspected for this draft.
\item Supplied literature source: \emph{Continuous time limit of the stochastic ensemble Kalman inversion}. \url{https://doi.org/10.1137/21M1437561}.
\item Supplied literature source: \emph{A Strongly Convergent Numerical Scheme from Ensemble Kalman Inversion}. \url{https://doi.org/10.1137/17M1132367}.
\item Literature assessment dated exactly \texttt{2026-08-21}: the supplied curation labels the problem partially solved and reports convergence in probability to continuous-time paths for nonlinear forward maps, convergence in moments for linear maps, and an unresolved strong-convergence question in the general nonlinear setting. These are supplied dated curation claims, not independently verified mathematical or current-literature conclusions. The supplied background explicitly requests independent verification of the time-sensitive status and cited claims.
\end{itemize}
\end{document}
